\begin{center}
\thispagestyle{empty}
\vspace{2cm}
\LARGE{\textbf{RESUMEN}}\\[1.0cm]
\end{center}
\thispagestyle{empty}

Este informe presenta el avance del proyecto “Lector PDF Minimalista”, una aplicación de escritorio orientada a la lectura y gestión de documentos PDF. El proyecto surge como respuesta a la necesidad de contar con una herramienta sencilla, organizada y centrada en el contenido, evitando la sobrecarga de botones, menús y funciones que caracteriza a algunos lectores disponibles actualmente. La aplicación está siendo desarrollada principalmente para Windows utilizando Electron, JavaScript, HTML, CSS y la biblioteca PDF.js.

Durante el primer sprint se construyó el prototipo inicial de la aplicación y se establecieron las bases de su arquitectura. Esta primera entrega permite abrir y visualizar archivos PDF, desplazarse entre sus páginas, consultar el número de página actual, ajustar el nivel de zoom, buscar y copiar texto, visualizar la estructura del documento mediante el outline y utilizar un modo de presentación. También se incorporaron funciones básicas de gestión, como eliminar páginas, guardar el documento y combinar dos archivos PDF en un solo documento, incluyendo una vista previa del resultado antes de guardarlo.

Además, el informe documenta los requerimientos iniciales del sistema, el producto mínimo viable correspondiente al Sprint 1, la organización del equipo, la arquitectura propuesta, los prototipos de interfaz y el proceso de control de calidad. Esta entrega representa la base funcional sobre la cual se desarrollarán las siguientes etapas del proyecto, que incluirán las anotaciones, los comentarios, la firma gráfica y la integración con el sistema de firma digital GAUDI, así como el soporte para formatos EPUB, DVI y Markdown. \\

\textbf{Palabras clave: }lector PDF, Electron, PDF.js, gestión de documentos, interfaz minimalista, combinación de PDF, aplicación de escritorio, Windows, GAUDI.